% 基礎数学 第11回 ワークシート（CC BY-SA 4.0）
% 解答版は worksheet_answers.tex をビルドする
\documentclass[a4paper,11pt]{ltjsarticle}
\usepackage{lfs-worksheet}

\begin{document}
\lfsheader{基礎数学}{第11回 行列の基礎 ワークシート}

\problem{$A = \begin{pmatrix} 1 & 2 \\ 3 & 4 \end{pmatrix}$，\ $B = \begin{pmatrix} 2 & 0 \\ -1 & 1 \end{pmatrix}$ のとき、次を計算しなさい。}
\begin{minipage}[t]{0.48\linewidth}
  (1) $A + B$
  \answerbox[2.6cm]{$\begin{pmatrix} 3 & 2 \\ 2 & 5 \end{pmatrix}$}
  (3) $AB$
  \answerbox[3.4cm]{$\begin{pmatrix} 1 \cdot 2 + 2 \cdot (-1) & 1 \cdot 0 + 2 \cdot 1 \\ 3 \cdot 2 + 4 \cdot (-1) & 3 \cdot 0 + 4 \cdot 1 \end{pmatrix} = \begin{pmatrix} 0 & 2 \\ 2 & 4 \end{pmatrix}$}
\end{minipage}\hfill
\begin{minipage}[t]{0.48\linewidth}
  (2) $2A - B$
  \answerbox[2.6cm]{$\begin{pmatrix} 2 & 4 \\ 6 & 8 \end{pmatrix} - \begin{pmatrix} 2 & 0 \\ -1 & 1 \end{pmatrix} = \begin{pmatrix} 0 & 4 \\ 7 & 7 \end{pmatrix}$}
  (4) $BA$
  \answerbox[3.4cm]{$\begin{pmatrix} 2 \cdot 1 + 0 \cdot 3 & 2 \cdot 2 + 0 \cdot 4 \\ -1 \cdot 1 + 1 \cdot 3 & -1 \cdot 2 + 1 \cdot 4 \end{pmatrix} = \begin{pmatrix} 2 & 4 \\ 2 & 2 \end{pmatrix}$}
\end{minipage}

\problem{次の積を計算しなさい。}
\[
  \begin{pmatrix} 1 & 0 & 2 \\ 3 & -1 & 1 \end{pmatrix}
  \begin{pmatrix} 1 \\ 2 \\ 3 \end{pmatrix}
\]
\answerbox[2.6cm]{$\begin{pmatrix} 1 + 0 + 6 \\ 3 - 2 + 3 \end{pmatrix} = \begin{pmatrix} 7 \\ 4 \end{pmatrix}$}

\problem{次の積は計算できるか。できる場合は、結果が何行何列の行列になるか答えなさい。}
\begin{minipage}[t]{0.48\linewidth}
  (1) $(2 \times 3$ 行列$) \times (3 \times 4$ 行列$)$
  \answerbox[2cm]{計算できる。$2 \times 4$ 行列}
\end{minipage}\hfill
\begin{minipage}[t]{0.48\linewidth}
  (2) $(3 \times 2$ 行列$) \times (3 \times 2$ 行列$)$
  \answerbox[2cm]{計算できない（左の列の数 $2$ と右の行の数 $3$ が違う）}
\end{minipage}

\problem{点 $(3,\, 1)$ を、原点のまわりに $90^\circ$ 回転させた点の座標を、行列 $\begin{pmatrix} 0 & -1 \\ 1 & 0 \end{pmatrix}$ を使って求めなさい。}
\answerbox[3cm]{$\begin{pmatrix} 0 & -1 \\ 1 & 0 \end{pmatrix}\begin{pmatrix} 3 \\ 1 \end{pmatrix} = \begin{pmatrix} -1 \\ 3 \end{pmatrix}$ \quad より \ $(-1,\, 3)$}

\problem{$A = \begin{pmatrix} 1 & 2 \\ 3 & 4 \end{pmatrix}$ と単位行列 $E$ について、$AE$ を計算し、$A$ と等しくなることを確かめなさい。}
\answerbox[3cm]{$\begin{pmatrix} 1 \cdot 1 + 2 \cdot 0 & 1 \cdot 0 + 2 \cdot 1 \\ 3 \cdot 1 + 4 \cdot 0 & 3 \cdot 0 + 4 \cdot 1 \end{pmatrix} = \begin{pmatrix} 1 & 2 \\ 3 & 4 \end{pmatrix} = A$}

\end{document}
